% =====================================================================================================
\chapter{Limitations}\label{ch:limitations}
% =====================================================================================================
The phase chapters list the caveats of each version. This chapter collects those that limit the project as a whole, ordered by
how much they constrain the conclusions.

\section{Sample size and labels}
\begin{itemize}
\item \textbf{Sixteen dynamically confirmed BHs.} Across RXTE and NICER only 16 dynamically confirmed BHs have usable data, and 7--9 of
them per instrument have hard-state observations. One source changes a source-level balanced accuracy by $\approx0.07$
(e.g.\ XTE J1118+480), and every bootstrap interval is wide because it resamples these few sources. The remaining BHs of
\citet{Marcel2026} have no usable RXTE or NICER timing data.
\item \textbf{Label definitions.} BHs are dynamically confirmed (BlackCAT, then \citealt{Marcel2026}); NSs are type-I bursters
(\citealt{Galloway2008}, MINBAR) and pulsars. The NS class mixes atolls, Z sources, AMXPs and globular-cluster sources, whose hard
states differ; BH candidates are used only in sensitivity analyses (v3c). The v2 source rule was executed incompletely (13 NSs missed,
found in v4b2), so the v1--v3 numbers describe a smaller NS sample than the rule implies.
\item \textbf{Population differences between samples.} Training BHs are LMXB transients; external sets add HMXBs (Cyg X-1, LMC X-1, LMC X-3)
and SS 433, external NSs are mostly AMXPs; this complicates the external comparisons in either direction.
\end{itemize}

\section{Independence and pre-registration}
\begin{itemize}
\item \textbf{Every pre-registration was written after the results of all earlier versions}, and from v5 onwards each was designed and
approved by the analyst under a standing instruction of the user, without external review. Later hypotheses were therefore guided by
earlier data (\cref{tab:synthladder}).
\item \textbf{The strongest tests are observation-level.} v11 and v13 used observations never downloaded before, with frozen pipelines and
rules, but of the same sources. v10 pooled all data but chose its hypotheses and pooling after the components. No test is
source-level independent.
\item \textbf{Post hoc elements} are labelled throughout (e.g.\ the source-equal-weight scale fix, the MINBAR matching rule, the v8c estimator
change and null-band subtraction, the v12 choice of RF), and none of them is used as a primary conclusion except the v12 RF choice, which
v13 then tested with RF pre-chosen.
\item \textbf{Multiplicity.} Descriptive comparisons were many (76 in v4a, 45 in v5, 16+16 in v4b1, dozens in v7--v9) and uncorrected; they are
reported but not interpreted individually.
\end{itemize}

\section{Measurement}
\begin{itemize}
\item \textbf{Count space.} RXTE spectral features are net count rates on a fixed channel-energy grid, not response-unfolded fluxes; the
power-law-normalised R and H$_R$ only partly remove response differences; PCU-combination tilts of $\approx\pm10\%$ at 25\,keV remain.
\item \textbf{Backgrounds.} RXTE uses the Standard-Products background model without systematic errors; NICER timing is not background-subtracted
before v12; 3C50 models only the mean background, uses its own GTI filtering, cannot process data with the original gain calibration
(139 + 38 faint observations excluded) and does not cover data after May 2023.
\item \textbf{Timing.} Standard-1 has no energy resolution and stops at 4\,Hz; the RXTE state index therefore omits 4--10\,Hz and uses BH
thresholds for NSs; event-mode estimates failed their consistency gate (v8c); NICER v9 used only file prefixes. The two $\nu_c$
definitions (Standard-1 and NICER) differ in band and bin weighting, so their $D$ values are not strictly comparable.
\item \textbf{States} are operational variability classes (rms thresholds), not full spectral-state assignments.
\item \textbf{No physical normalisation.} Eddington ratio, distance, mass and inclination are not controlled; brightness enters only as a count
rate. The $\nu_c$ result cannot be attributed to mass rather than accretion state.
\end{itemize}

\section{Statistics}
\begin{itemize}
\item \textbf{Fixed-prediction resampling.} All bootstrap intervals and permutation tests keep the out-of-fold predictions or features fixed; they
do not include the variability of model fitting or of nested model selection. Nested source AUCs are approximations (shifted scores).
\item \textbf{Observation-level metrics within states} pool the observations of the drawn sources and are dominated by sources with many
observations; the source-level versions are reported alongside but have even less power.
\item \textbf{Saturation.} Headline metrics saturated on the external set of v6 (AUC 1.000), which the decision rule did not anticipate.
\item \textbf{The colour-conditional permutation test} residualises with a label-free kNN; when BHs cluster in colour they become each other's
neighbours, which makes the test conservative.
\end{itemize}

\section{Process and reproducibility}
\begin{itemize}
\item Visual verdicts (burst candidates, figure checks) were made by one person.
\item Raw data are not in the repository (by design); they can be re-downloaded with the scripts, and every file's URL, size and SHA256 is in
\file{logs/downloads.jsonl}. Some NICER raw bytes were streamed and not stored (v11).
\item Analyses ran on Windows (Python 3.12.10) and, for some re-checks, macOS (Python 3.12.14); feature files reproduced to $\le10^{-16}$ and RF source
means to $\le0.001$ across machines. HEASoft steps (XSPEC, \code{nh}, \code{nibackgen3C50}) ran in WSL.
\item This report re-ran no analysis; all numbers are read from committed files, and a small number of descriptive quantities (per-source
hard-like means in Phase 4, the v9b Mann--Whitney check) were recomputed read-only from committed CSVs, as stated where they appear.
\end{itemize}
